\scanpage{11}
\deleted{\illegible{1}} Les pts \inserted{$x$} de $\underline{\pi_0}(X/k)$ correspondent aux \uline{composantes} \uline{connexes} de $X$, pour toute telle composante $X_{(x)}$, $k(x)$ est la plus grande sous-algèbre entière séparable de $\Gamma(X_{(x)},\underline{O}_{X_{(x)}})$ sur $k(x)$.

Que se passe-t-il pour $X=\operatorname{Spec}k'$ ? \deleted{\unsure{de|du}}

On a $k'=\varinjlim k_i$,\inserted{(\unsure{sous-ext.|sous-cat.} \unsure{galoisiennes|galoisienne} \unsure{finies|finie} de $k'$)} $\Rightarrow X=\varprojlim X_i$, $X_i=\operatorname{Spec} k_i$.

\[
\begin{aligned}
X'&=\varprojlim X'_i=\varprojlim_i(\pi_i)_{k}=(\pi)_k\\
\pi_0(X')&=\pi\quad(\text{isom. \uline{topologique}})
\end{aligned}
\]
\deleted{Ce n'est plus un schéma étale sur $k$}

Schéma proétale sur $k$ [\deleted{compact, réduit, à corps résiduels} \inserted{$=$ entier séparable $=$} \deleted{\unsure{à points ext.|à points isolés}} réduit à ext. résiduelles \deleted{\unsure{étal|étal.}} alg. sép.]

\deleted{\S} [Schéma loc. proétale sur $k$]

\[
\left\{\begin{aligned}
\mathrm{ProEt}(k)&\cong\pi\text{-esp cpt tdisc.}\cong\mathrm{Alg}_{\text{\unsure{indét|indet|ind-et}}}(k)\\
\mathrm{Locproet}(k)&\cong\pi\text{-esp loc. cpt disc.}
\end{aligned}\right.
\]
Si $X$ qu. cpct sur $k$ (pour simplifier) on construit une algèbre \unsure{indétale|proétale} \deleted{sur $X$} $k$, au choix

a)

\[
\begin{aligned}
A&=\text{limite inductive des sous-}\underline{\unsure{\mathrm{Alg}^{\prime\prime}\mathrm{bres}|\mathrm{Al}''\mathrm{bres}}}\ \text{étales de }\Gamma(X,\underline{O}_X)\\
&=\text{\deleted{limite} plus grande sous-Alg. ind-étale de }\Gamma(X,\underline{O}_X).\\
X&\longrightarrow\operatorname{Spec}(A)\overset{\text{dfn}}{=}\underline{\pi}_0(X/k)
\end{aligned}
\]
b) Pour toute ext. \deleted{finie} \inserted{étale} $k_i\subset k'$ de $k$, on désigne par \deleted{\unsure{$k$|$A$}} $A_i$ l'anneau des fonctions loc. constantes sur $X_{k_i}$, à valeurs dans $k_i$. Soit $\overline A=\varinjlim_i {A_i}_{k'}$, \inserted{fonctions loc. constantes sur $X_{k_i}$ \unsure{à v.|à val} dans $k'$} \deleted{\unsure{alors|ainsi} $\pi$ opère} \inserted{fonctions} loc. constantes sur $X_{k'}$ à valeurs dans $k'$, Alors \inserted{$\overline A$ est ind-étale et} $\pi$ opère dessus par transfert des structures, on pose $\overset{\deleted{\unsure{\text{donc}|\text{alors}}}}{A}=(\overline A)^\pi$.

\par\noindent\rule{.20\linewidth}{.4pt}\par
\[
\left\{\begin{array}{ll}
\text{a)} & \parbox[t]{.82\linewidth}{$\deleted{X\to}\underline{\pi}_0(X/k)$ est proétale et $X\to\underline{\pi}_0(X/k)$ \deleted{est} surj. à fibres géom. connexes}\\[.7ex]
\text{b)} & \underline{\pi}_0(X/k)\ \text{fonctoriel en }X\\[.7ex]
\text{c)} & \underline{\pi}_0(X/k)\ \text{commute au produit}\\[.7ex]
\text{d)} & \underline{\pi}_0(X/k)\ \text{commute à l'extension de la base ....}
\end{array}\right.
\]
\uline{NB} $\underline{\pi}_0(X/k)$ peut se définir comme schém. loc proétale, \deleted{si} toute pt de \deleted{ce} $X$ a un voisinage ouvert et fermé $U$, tel que \deleted{\unsure{la|le}} la décomposition de $U_{k'}$ en somme \deleted{\unsure{de|s}} \deleted{soit} soit finie (p. ex. $U$ qu. cpct).

% archive sheet number: 10
